% Private notes, not part of the submission.  Written for the October draft: section, table and notebook
% references (SM Sec. S*, Tables S*, notebook 03) follow that draft and its sm.tex, both removed (git
% history); the derivations themselves still hold.
% Companion to paper.tex and sm.tex: every derivation written out step by step, with the
% assumptions of each step, independent numerical checks, and the points that need attention.
% Build: latexmk -pdf derivations.tex   (from this directory)
\documentclass[11pt]{article}
\usepackage[margin=2.6cm]{geometry}
\usepackage{amsmath,amssymb,amsthm,bm}
\usepackage{xcolor}
\usepackage[hidelinks]{hyperref}
\usepackage{booktabs}

\newcommand{\E}{\mathbb E}
\newcommand{\Var}{\operatorname{Var}}
\newcommand{\Cov}{\operatorname{Cov}}
\newcommand{\Sb}{\bar S}
\newcommand{\PhiN}{\Phi_{\mathcal N}}
\newcommand{\fb}{\bar f}
\newcommand{\dd}{\mathrm d}
\newcommand{\exact}{\textsf{\small[exact]}}
\newcommand{\controlled}{\textsf{\small[controlled]}}
\newcommand{\heur}{\textsf{\small[heuristic]}}
\newcommand{\numer}{\textsf{\small[numerical]}}

\newtheorem{proposition}{Proposition}
\newenvironment{verify}{\par\medskip\noindent\textcolor{teal}{\textbf{Check.}}\ }{\par\medskip}
\newenvironment{flag}{\par\medskip\noindent\textcolor{red!70!black}{\textbf{Flag.}}\ }{\par\medskip}

\begin{document}

\title{Derivations behind ``Interaction heterogeneity and the multiscaling of firm growth volatility''}
\author{Ludovico Furlanetto}
\date{Working note, \today}
\maketitle

\begin{abstract}
This note writes out every calculation of the article and its Supplemental Material, and of the
deterministic life-cycle theory of \texttt{notebooks/03\_life\_cycles.ipynb}, one step at a time, with
the assumption each step uses.  Each result carries a tag: \exact{} (follows from the model with no
approximation), \controlled{} (leading order in a stated small parameter), \heur{} (a physically motivated
approximation whose error is not bounded), or \numer{} (established only by simulation).  Where a closed
form could be tested independently of the notebooks it was, and the test is reported in a \emph{Check}
paragraph.  Places where the derivation needs a caveat, a correction or more care than the article gives
it are marked \emph{Flag}, and Sec.~\ref{sec:flags} collects them.
\end{abstract}

\tableofcontents

%======================================================================
\section{The model and notation}
\label{sec:model}

Firms $i=1,\dots,N$ sit on the nodes of a configuration-model graph with adjacency matrix $A$, degree law
$p(k)\propto k^{-\alpha}$ and mean degree $C$.  Sizes $x_i>0$ obey
\begin{equation}
\dd x_i=x_i\Big[1-\frac{x_i}{m}+\sum_jJ_{ij}\frac{x_j}{m}\Big]\dd t+\sqrt{D\,m\,x_i}\;\dd W_i,\qquad
m=\frac1N\sum_jx_j,
\label{eq:model}
\end{equation}
with independent Wiener processes $W_i$ and
\begin{equation}
J_{ij}=A_{ij}\Big(\frac{\mu}{C}+\frac{\sigma z_{ij}}{\sqrt C}\Big),\qquad z_{ij}\sim\mathcal N(0,1)\ \text{i.i.d.},\quad
\E[z_{ij}z_{ji}]=0 .
\label{eq:J}
\end{equation}
The reciprocity is zero.  $\mu<0$ is net competition (the code uses the opposite sign: \texttt{mu=1.76}
means $\mu=-1.76$).  When $x_i$ reaches $0$ the firm exits and a newborn of size $Dm$ takes node $i$ with
its couplings.

\paragraph{Notation.}  $S_i=x_i/m$ is the relative size; $\kappa_i=k_i/C$ the relative degree, so that the
degree average $\langle\kappa\rangle=1$ by construction, and $V\equiv\Var\kappa=\langle\kappa^2\rangle-1$.
$\langle\cdot\rangle$ is an average over degree classes (or over nodes), $\E$ an average over the Gaussian
disorder or the noise, and an overbar a time or life average.  $\phi$ and $\PhiN$ are the standard normal
density and distribution function, $\phi_0\equiv\phi(0)=1/\sqrt{2\pi}$.

\paragraph{Homogeneity.}  Every term of Eq.~\eqref{eq:model} is homogeneous of degree one in the vector $x$:
under $x\to cx$, $m\to cm$, the drift scales by $c$ and the noise amplitude $\sqrt{Dmx_i}$ also by $c$.  So
if $x(t)$ is a solution, $cx(t)$ is a solution with $W$ unchanged.  The economy has no preferred size, and
only relative sizes carry information.  \exact

%======================================================================
\section{Relative sizes: the one-firm equation}
\label{sec:relative}

\subsection{The mean firm}

Write $F_i=1-S_i+\sum_jJ_{ij}S_j$, so that the drift of $x_i$ is $x_iF_i$ (because $x_j/m=S_j$).  Summing
Eq.~\eqref{eq:model} over $i$ and dividing by $N$,
\begin{equation}
\dd m=\frac1N\sum_jmS_jF_j\,\dd t+\frac1N\sum_j\sqrt{D\,m\cdot mS_j}\,\dd W_j
=m\big[\fb_0\,\dd t+\dd Z\big],
\label{eq:dm}
\end{equation}
with
\begin{equation}
\fb_0=\frac1N\sum_jS_jF_j,\qquad \dd Z=\frac1N\sum_j\sqrt{DS_j}\,\dd W_j,\qquad
(\dd Z)^2=\frac{1}{N^2}\sum_jDS_j\,\dd t=\frac DN\,\dd t ,
\label{eq:Z}
\end{equation}
where the last equality uses $\sum_jS_j=N$.  The common noise $Z$ has quadratic variation of order
$1/N$.  \exact

\subsection{It\^o's rule for the quotient}

For $S_i=x_i/m$ It\^o's rule for $f(x,m)=x/m$ gives
\begin{equation}
\dd S_i=\frac{\dd x_i}{m}-\frac{x_i\,\dd m}{m^2}+\frac{x_i\,(\dd m)^2}{m^3}-\frac{\dd x_i\,\dd m}{m^2}.
\end{equation}
Term by term:
\begin{align}
\frac{\dd x_i}{m}&=S_iF_i\,\dd t+\sqrt{DS_i}\,\dd W_i,\\
\frac{x_i\,\dd m}{m^2}&=S_i\big[\fb_0\,\dd t+\dd Z\big],\\
\frac{x_i(\dd m)^2}{m^3}&=S_i\,(\dd Z)^2=\frac{DS_i}{N}\,\dd t,\\
\frac{\dd x_i\,\dd m}{m^2}&=\sqrt{DS_i}\,\dd W_i\cdot\frac1N\sqrt{DS_i}\,\dd W_i=\frac{DS_i}{N}\,\dd t ,
\end{align}
where the last line uses $\dd W_i\dd W_j=\delta_{ij}\dd t$.  The two It\^o corrections cancel exactly, and
\begin{equation}
\boxed{\dd S_i=S_i\big[F_i-\fb_0\big]\dd t+\sqrt{DS_i}\,\dd W_i-S_i\,\dd Z .}
\label{eq:rel}
\end{equation}
\exact{} (between entries).

\begin{flag}
SM Eq.~(S1) writes a remainder ``$+O(N^{-1})\,\dd t$''.  The remainder is exactly zero: the
$x(\dd m)^2/m^3$ and $\dd x\,\dd m/m^2$ terms cancel identically.  Harmless, but the SM can say so.
\end{flag}

\paragraph{Conservation.}  Summing Eq.~\eqref{eq:rel},
$\sum_i\dd S_i=\big(\sum_iS_iF_i-N\fb_0\big)\dd t+\sum_i\sqrt{DS_i}\,\dd W_i-N\,\dd Z=0$
by the definitions of $\fb_0$ and $Z$.  So $\sum_iS_i=N$ is preserved between entries.  \exact

\paragraph{Size of the common-noise term.}  $S_i\,\dd Z$ has variance $S_i^2D\,\dd t/N$, against $DS_i\,\dd t$
for the own noise: relative size $S_i/N$.  It is of order $N^{-1/2}$ in amplitude and is dropped in the
one-firm equation.  It is the only place where the finite size of the economy enters the single-firm
dynamics directly.  \controlled{} in $1/N$.

\subsection{Entry dilution}

An exit at $S_i=0$ removes nothing from $m$.  The newborn adds $Dm$ to $\sum_jx_j$, so
$m\to m'=m(1+D/N)$, every incumbent's relative size is multiplied by $1/(1+D/N)$, and the newborn's
relative size is $D/(1+D/N)=D+O(1/N)$.  If entries happen at a rate $N\nu$ per unit time ($\nu$ per
node), then over $\dd t$
\begin{equation}
\frac{m(t+\dd t)}{m(t)}=(1+D/N)^{N\nu\,\dd t}\big[1+\fb_0\,\dd t+\dd Z\big]
\;\Rightarrow\; \fb=\fb_0+D\nu ,
\end{equation}
and each incumbent picks up the drift $-D\nu S_i$.  So the growth rate of the mean firm is
\begin{equation}
\fb(t)=\frac1N\sum_iS_iF_i+D\nu(t),
\label{eq:fbar}
\end{equation}
and in Eq.~\eqref{eq:rel} $\fb_0$ is replaced by $\fb$.  \exact{} at the level of rates; $\nu(t)$ is
the empirical entry rate, which fluctuates at finite $N$.  (At the reference point $D\nu\simeq0.06$.)

\subsection{The field and the one-firm equation}

Collect everything except the firm's own self-limitation into the field
\begin{equation}
\Phi_i(t)=F_i+S_i-\fb=1+\sum_jJ_{ij}S_j(t)-\fb(t).
\label{eq:field}
\end{equation}
Then
\begin{equation}
\boxed{\dd S_i=S_i\big[\Phi_i(t)-S_i\big]\dd t+\sqrt{DS_i}\,\dd W_i+O(N^{-1/2}),}
\label{eq:onefirm}
\end{equation}
with exit at $S_i=0$ and entry at $S_i=D$.  This is exact as a rewriting; the content of the
``one firm is enough'' result is that $\Phi_i$ may be treated as an external input, i.e.\ that firm $i$'s own
noise does not feed back on $\Phi_i$ through its neighbours.  The feedback is of order
$J_{ij}J_{ji}\sim A_{ij}(\mu^2/C^2+\sigma^2\E[z_{ij}z_{ji}]/C)$, so at zero reciprocity it is $O(1/C)$ per
neighbour and $O(1/C)$ in total after summing $k\sim C$ neighbours with weight $1/C$ each.  \controlled{}
in $1/C$, tested \numer{} (SM Table~S3).

\paragraph{Plateau.}  At constant $\Phi>0$ the noise-free equation $\dot S=S(\Phi-S)$ has the stable fixed
point $S=\Phi$ ($\partial_S[S(\Phi-S)]=\Phi-2S=-\Phi<0$ there) and relaxes to it at rate $\Phi$.  For
$\Phi<0$, $S=0$ is the only non-negative fixed point and is attracting.  In absolute units the plateau is
$x_i=m\Phi$.

\subsection{Logarithmic form and ``size is mean field''}

By It\^o, $\dd\ln S=\dd S/S-\tfrac12(\dd S)^2/S^2$ with $(\dd S)^2=DS\,\dd t$:
\begin{equation}
\dd\ln S_i=\Big[\Phi_i-S_i-\frac{D}{2S_i}\Big]\dd t+\sqrt{\frac{D}{S_i}}\,\dd W_i .
\label{eq:log}
\end{equation}
Integrate over a window $[t_1,t_2]$ of length $T$ with $\ln S_i(t_2)=\ln S_i(t_1)$ and divide by $T$:
\begin{equation}
\Sb_i=\bar\Phi_i-\frac D2\,\overline{S_i^{-1}}+\frac1T\int_{t_1}^{t_2}\sqrt{D/S_i}\,\dd W_i .
\end{equation}
Without noise, $\Sb_i=\bar\Phi_i$ exactly: a firm's mean size is its mean field.  With noise the
stochastic integral has zero mean for a window fixed in advance, but not for a window chosen by the path
(a life selected for its length), which is why the identity holds per life only to about 5\%.  \exact{}
(the first equality); the per-life accuracy is \numer.

\subsection{Mapping to the standard GLV with demographic noise}

Suppose $\fb(t)\to f^*$ (self-averaging at large $N$) and $r\equiv1-f^*>0$.  Then
$\dd S_i=S_i[r-S_i+\sum_jJ_{ij}S_j]\dd t+\sqrt{DS_i}\,\dd W_i$.  Set $S=ry$, $t'=rt$.
\begin{itemize}
\item Drift: $S[r-S+\sum JS]\dd t=ry\cdot r[1-y+\sum Jy]\cdot\dd t'/r=ry[1-y+\sum Jy]\,\dd t'$.
\item Noise: $W_t$ as a function of $t'$ has $\E[(\dd W)^2]=\dd t=\dd t'/r$, so $\dd W=r^{-1/2}\dd W'$ with $W'$
  standard in $t'$.  Then $\sqrt{DS}\,\dd W=\sqrt{Dry}\,r^{-1/2}\dd W'=\sqrt{Dy}\,\dd W'$.
\item Divide by $r$ (since $\dd S=r\,\dd y$):
\end{itemize}
\begin{equation}
\dd y_i=y_i\Big[1-y_i+\sum_jJ_{ij}y_j\Big]\dd t'+\sqrt{\frac{D}{r^2}\,y_i}\;\dd W_i' ,
\end{equation}
with newborns at $y=D/r$.  This is the GLV with demographic noise of strength $D'=D/r^2$.
\controlled{} (requires $\fb$ to self-average and $f^*<1$).

\begin{flag}
In the $y$ variables the newborn size is $D/r=D'r$, not $D'$: the identity ``newborn size $=$ noise unit''
of the original model survives the mapping only at zero growth ($r=1$).  The SM states the newborn size
correctly as $D/r$; when comparing with Ref.~\cite{altieri2021} it is worth saying that the entry size
is a second renormalised parameter.
\end{flag}

\paragraph{The deterministic (floor) version.}  Replace the noise by $\lambda(m-x_i)$.  In relative units this
adds $\lambda(1-S_i)$ to $\dot S_i$, and adds $N^{-1}\sum_j\lambda(1-S_j)=0$ to $\fb$: the floor does not shift
the growth rate directly.  In the $y$ variables, $\lambda(1-ry)\,\dd t/r=(\lambda/r^2-\lambda y/r)\,\dd t'$: an
immigration $\lambda/r^2$ and a shift $-\lambda/r$ of the intrinsic growth rate.  In logs,
$\dd\ln S/\dd t=\Phi-S+\lambda(S^{-1}-1)$, the starting point of Sec.~\ref{sec:lifecycle}.  \exact

%======================================================================
\section{Protocol and estimators}
\label{sec:protocol}

\paragraph{Lives.}  Sizes are recorded every $\Delta t=0.5$ (``one year'').  A life is a maximal run of
samples at one node with the same death counter; lives cut by the window are kept.  A life with $n+1$
samples has $n$ growth rates $g_t=\ln S(t+\Delta t)-\ln S(t)$.  Only lives with $n\ge n_{\min}$ enter.

\paragraph{Volatility.}  $v=\sqrt{\pi/2}\cdot\frac1n\sum_t|g_t-\bar g|$.  The factor: if
$X\sim\mathcal N(\mu,s^2)$, then $\E|X-\mu|=s\int|z|\phi(z)\,\dd z=2s\int_0^\infty z\phi(z)\,\dd z=2s\phi_0=s\sqrt{2/\pi}$,
so $\sqrt{\pi/2}\,\E|X-\mu|=s$.  For Gaussian growth $v$ estimates the standard deviation; for fat-tailed
growth it is a robust scale that is smaller than the standard deviation.  \exact

\paragraph{Exponents.}  Lives are sorted by $\Sb$ (the life-mean relative size) into 25 bins of equal count;
in each bin $\E[v^q\mid\Sb]$ is the bin mean of $v^q$ and $\Sb$ the bin mean of $\Sb$; $-\zeta_q$ is the
least-squares slope of $\ln\E[v^q\mid\Sb]$ on $\ln\Sb$ over all bins.  $\zeta_q/\zeta_1$ is the ratio of the
fitted slopes.  This is a definition; the local slopes ``by fifths'' are the same fit restricted to five
bins at a time.

\paragraph{Why $\zeta_q=q\zeta_1$ for a fixed-shape law.}  If $v=A(\Sb)X$ with the law of $X$ independent of
$\Sb$, then $\E[v^q\mid\Sb]=A(\Sb)^q\E[X^q]$ and $\zeta_q=-q\,\dd\ln A/\dd\ln\Sb=q\zeta_1$, for every $q$ at
every size.  \exact{}  The general statement is Sec.~\ref{sec:cumulant}.

%======================================================================
\section{Heterogeneous dynamical mean-field theory}
\label{sec:hdmft}

\subsection{Degree classes and kernels}

Group the nodes into classes $p$ with population fractions $\rho_p$ and degrees $k_p$.  Let $P_{pr}$ be the
probability that a given node of class $p$ is linked to a given node of class $r$.  The expected number
of class-$r$ neighbours of a class-$p$ node is $N\rho_rP_{pr}$, and the degree constraint is
$\sum_rN\rho_rP_{pr}=k_p$.  Define
\begin{equation}
K_{pr}=\frac{N\rho_rP_{pr}}{C}\qquad\Rightarrow\qquad\sum_rK_{pr}=\kappa_p .
\end{equation}
\emph{Uncorrelated (bilinear) kernel:} $P_{pr}=k_pk_r/(NC)$, so $K_{pr}=\rho_r\kappa_p\kappa_r$; it satisfies
the constraint because $\sum_r\rho_r\kappa_r=1$.  It is a probability only if $k_pk_r\le NC$; at $\alpha=2.5$
the largest degrees violate this (structural cutoff $k\sim\sqrt{NC}$).  \emph{Soft configuration:}
$P_{pr}=h_ph_r/(h_ph_r+C/N)$, always in $[0,1]$, with hidden degrees $h_p$ chosen so that the constraint
holds; for $h_ph_r\ll C/N$ it reduces to the bilinear form with $h_p\propto k_p$.

\subsection{The mean part of the interaction field}

The interaction term on node $i$ in class $p$ is
$\sum_jJ_{ij}S_j=\frac\mu C\sum_jA_{ij}S_j+\frac{\sigma}{\sqrt C}\sum_jA_{ij}z_{ij}S_j$.
The first sum has $k_i\sim C$ terms; by the law of large numbers over neighbours drawn with the kernel,
\begin{equation}
\frac\mu C\sum_jA_{ij}S_j\;\to\;\frac\mu C\sum_rN\rho_rP_{pr}\,\E[S_r]=\mu\sum_rK_{pr}\E[S_r]\equiv\mu B_p ,
\end{equation}
with fluctuations of relative order $k^{-1/2}$, i.e.\ an absolute correction $O(C^{-1/2})$.
\controlled{} in $1/C$.

\subsection{The random part: cavity argument}

Remove node $0$ (the cavity) and let the rest of the system evolve; the neighbours' sizes $S_j^{\setminus0}(t)$
are then independent of the couplings $z_{0j}$.  Reinserting node $0$ has two effects.
\begin{enumerate}
\item \emph{Direct field.}  $\eta_0(t)=C^{-1/2}\sum_jA_{0j}z_{0j}S_j^{\setminus0}(t)$ is a sum of $k_0$
  independent zero-mean terms conditional on the cavity trajectories, so by the central limit theorem it is a
  Gaussian process with covariance
  \begin{equation}
  \E[\eta_0(t)\eta_0(t')]=\frac1C\sum_jA_{0j}S_j(t)S_j(t')\;\to\;\sum_rK_{pr}\E[S_r(t)S_r(t')]\equiv q_p(t,t') .
  \end{equation}
\item \emph{Reaction.}  Node $0$ perturbs each neighbour $j$ by $J_{j0}S_0$; the response of $j$ is
  $\chi_j J_{j0}S_0$ and feeds back on $0$ through $J_{0j}$.  The total is
  $\sum_jJ_{0j}J_{j0}\chi_jS_0$, whose mean is $\sum_jA_{0j}(\mu^2/C^2+\sigma^2\E[z_{0j}z_{j0}]/C)\chi_j=O(\mu^2/C)+0$
  at zero reciprocity.  It vanishes for $C\to\infty$.
\end{enumerate}
So for $C\to\infty$ the field on a class-$p$ node is
\begin{equation}
\Phi_p(t)=1-\fb(t)+\mu B_p(t)+\sigma\eta_p(t),\qquad \E[\eta_p(t)\eta_p(t')]=q_p(t,t'),
\label{eq:effproc}
\end{equation}
\begin{equation}
B_p(t)=\sum_rK_{pr}\E[S_r(t)],\qquad q_p(t,t')=\sum_rK_{pr}\E[S_r(t)S_r(t')] .
\label{eq:closure}
\end{equation}
Each class obeys Eq.~\eqref{eq:onefirm} with this field; Eq.~\eqref{eq:closure} closes the problem, and
$\fb(t)$ is fixed by $\sum_p\rho_p\E[S_p(t)]=1$ (which is $\sum_iS_i=N$, not an extra approximation).
\controlled{} in $1/C$; this is the standard dynamical cavity derivation
\cite{bunin2017,galla2018,aguirrelopez} with degree classes.

\subsection{Bilinear kernel: the field law}

With $K_{pr}=\rho_r\kappa_p\kappa_r$, both closures factorise:
\begin{equation}
B_p=\kappa_p\bar B,\quad \bar B=\sum_r\rho_r\kappa_r\E[S_r];\qquad q_p=\kappa_pQ,\quad Q(t,t')=\sum_r\rho_r\kappa_r\E[S_r(t)S_r(t')] .
\end{equation}
$\bar B$ and $Q$ are degree-weighted averages (a random neighbour is drawn in proportion to its degree).
Writing $\eta_p=\sqrt{\kappa_p}\,\xi$ with $\E[\xi(t)\xi(t')]=Q(t,t')$:
\begin{equation}
\boxed{\Phi_\kappa(t)=1-\fb(t)+\mu\kappa\bar B+\sigma\sqrt\kappa\,\xi(t).}
\label{eq:fieldlaw}
\end{equation}
Under competition ($\mu<0$, $\bar B>0$) the field's mean falls linearly in $\kappa$; its amplitude grows as
$\sqrt\kappa$.  \controlled{}  (On the reference graph the bilinear kernel is not a probability for the hubs,
so the factorisation is approximate there; the measured amplitude exponent $0.54$ against $0.5$ is \numer.)

%======================================================================
\section{The deterministic static solution and its stability}
\label{sec:static}

\subsection{Stationary sizes at a static field}

Without noise and floor, if the field is static, $\Phi=h$, every firm goes to the attracting fixed point of
$\dot S=S(h-S)$: $S^*=h$ if $h>0$, $S^*=0$ if $h\le0$.  So $S^*=[h]_+$.  (For an extinct firm the condition
$h<0$ is also the non-invasibility condition.)  In class $p$, $h_p=a+\mu B_p^*+\sigma\sqrt{q_p^*}\,z$ with
$a=1-f^*$, $z\sim\mathcal N(0,1)$ frozen; so
\begin{equation}
S_p^*=\sigma\sqrt{q_p^*}\,[\Delta_p+z]_+,\qquad \Delta_p=\frac{a+\mu B_p^*}{\sigma\sqrt{q_p^*}} .
\end{equation}

\subsection{Truncated Gaussian moments}

Define $w_n(\Delta)=\int_{-\Delta}^\infty(\Delta+z)^n\phi(z)\,\dd z$.  Using $\int_{-\Delta}^\infty z\phi\,\dd z=\phi(\Delta)$ and
$\int_{-\Delta}^\infty z^2\phi\,\dd z=\PhiN(\Delta)-\Delta\phi(\Delta)$ (integration by parts with $z\phi=-\phi'$):
\begin{align}
w_0&=\PhiN(\Delta),\\
w_1&=\phi(\Delta)+\Delta\PhiN(\Delta),\\
w_2&=\Delta^2\PhiN+2\Delta\phi+\PhiN-\Delta\phi=(1+\Delta^2)\PhiN(\Delta)+\Delta\phi(\Delta).
\end{align}
Two identities are used repeatedly.
\begin{enumerate}
\item $w_n'(\Delta)=n\,w_{n-1}(\Delta)$ for $n\ge1$ (differentiate under the integral; the boundary term is
  $(\Delta+z)^n\phi|_{z=-\Delta}=0$), and $w_0'=\phi$.
\item $w_2-w_0=\Delta w_1$.  Proof: $(\Delta+z)^2=\Delta(\Delta+z)+z(\Delta+z)$ and
  $\int_{-\Delta}^\infty z(\Delta+z)\phi\,\dd z=\int(\Delta+z)(-\phi')\,\dd z=[-(\Delta+z)\phi]_{-\Delta}^\infty+\int_{-\Delta}^\infty\phi=w_0$.
\end{enumerate}
Values at $0$: $w_0(0)=\tfrac12$, $w_1(0)=\phi_0$, $w_2(0)=\tfrac12$; $w_0''(\Delta)=\phi'(\Delta)=-\Delta\phi(\Delta)$ vanishes at $0$.

\begin{verify}
Both identities verified numerically on a grid of $\Delta$ to $10^{-6}$.
\end{verify}

\subsection{Static closure}

The moments of $S_p^*$ are $\E S_p^*=\sigma\sqrt{q_p^*}\,w_1(\Delta_p)$, $\E(S_p^*)^2=\sigma^2q_p^*w_2(\Delta_p)$ and the
survival fraction $\varphi_p=\Pr(S_p^*>0)=w_0(\Delta_p)$.  Inserting into Eq.~\eqref{eq:closure} at equal times
($q_p^*$ is the static covariance):
\begin{equation}
B_p^*=\sum_rK_{pr}\,\sigma\sqrt{q_r^*}\,w_1(\Delta_r),\qquad
q_p^*=\sum_rK_{pr}\,\sigma^2q_r^*\,w_2(\Delta_r),\qquad
1=\sum_p\rho_p\,\sigma\sqrt{q_p^*}\,w_1(\Delta_p).
\label{eq:staticclosure}
\end{equation}
The last equation (normalisation) is what fixes $a$, i.e.\ the growth rate $f^*=1-a$.  \controlled

\subsection{Stability: the response argument}

Add to every node an independent, static, zero-mean perturbing field of variance $\epsilon$, and look for
the stationary response.  A surviving firm ($h>0$) responds with unit susceptibility, $\delta S=\delta h$; an
extinct one does not respond.  The perturbation of the neighbours' sizes creates, through the random
couplings, an extra field on a class-$p$ node of variance $\sigma^2\sum_rK_{pr}\E[\delta S_r^2]$ (same cavity
argument as above, with $\delta S$ in place of $S$).  In the cavity picture, whether a node survives is
independent of the new perturbation, so with $Y_p=\E[\delta S_p^2]$,
\begin{equation}
Y_p=\varphi_p\Big(\epsilon+\sigma^2\sum_rK_{pr}Y_r\Big)\quad\Longrightarrow\quad
Y=\big(\mathbb 1-\sigma^2\,{\rm diag}(\varphi)K\big)^{-1}\varphi\,\epsilon .
\end{equation}
The response is finite while the Perron root $\Lambda$ of $\sigma^2{\rm diag}(\varphi)K$ (a non-negative matrix) is
below one and diverges as $\Lambda\to1$: the unique fixed point loses stability there.  \controlled{} (this
is the standard criterion of Refs.~\cite{opperdiederich1992,bunin2017,roy2020}; the mean part of the
perturbation, $\mu\sum_rK_{pr}\E\delta S_r$, has zero mean and does not diverge).

\subsection{Regular and complete graphs: $\sigma_c=\sqrt2$}

With one class, $\kappa=1$, $K=1$, so $B^*=\E S^*=1$ by the normalisation, and the mean interaction enters
every field as the same constant $\mu$.  It can be absorbed into $a$: the composition does not depend on
$\mu$, only the growth rate does.  The closure reads
\begin{equation}
\sigma^2w_2(\Delta)=1,\qquad \sigma\sqrt{q^*}\,w_1(\Delta)=1,\qquad \Delta=\frac{a+\mu}{\sigma\sqrt{q^*}} .
\end{equation}
The first fixes $\Delta(\sigma)$ ($w_2$ is increasing, from $0$ to $\infty$, so the root is unique).  The second
gives $\sigma\sqrt{q^*}=1/w_1(\Delta)$, and the third $a+\mu=\Delta/w_1(\Delta)$, so
\begin{equation}
f^*=1-a=1+\mu-\frac{\Delta}{w_1(\Delta)}\equiv g_0(\sigma)+\mu,\qquad g_0(\sigma)=1-\frac{\Delta(\sigma)}{w_1(\Delta(\sigma))} .
\end{equation}
(The article writes $g_0-|\mu|$ for $\mu<0$.)  Stability: $\Lambda=\sigma^2\varphi=\sigma^2w_0(\Delta)$.  At the threshold
$\sigma^2w_0=1=\sigma^2w_2$, so $w_2-w_0=\Delta w_1=0$, and since $w_1>0$, $\Delta=0$.  Then $w_2(0)=\tfrac12$ gives
\begin{equation}
\boxed{\sigma_c=\sqrt2,\qquad \varphi_c=\tfrac12,}
\end{equation}
independent of $\mu$.  \controlled

\paragraph{Zero-growth line on the regular graph.}  $f^*=0$ gives $\mu^{(0)}(\sigma)=\Delta/w_1(\Delta)-1=-g_0(\sigma)$.  At
$\sigma=\sqrt2$, $\Delta=0$ and $\mu^{(0)}=-1$: the two regular-graph lines meet at $(\mu,\sigma)=(-1,\sqrt2)$.

\paragraph{Its slope at the meeting point.}  Parametrise by $\Delta$.  $\dd\mu/\dd\Delta=\dd(\Delta/w_1)/\dd\Delta=(w_1-\Delta w_0)/w_1^2$, which at
$\Delta=0$ is $1/\phi_0=\sqrt{2\pi}$.  From $\sigma=w_2^{-1/2}$, $\dd\sigma/\dd\Delta=-\tfrac12w_2^{-3/2}\cdot2w_1=-w_1w_2^{-3/2}$, which at $0$ is
$-\phi_0\cdot2^{3/2}=-2/\sqrt\pi$.  Hence
\begin{equation}
\frac{\dd\mu^{(0)}}{\dd\sigma}\Big|_{\sigma=\sqrt2}=\frac{\sqrt{2\pi}}{-2/\sqrt\pi}=-\frac{\pi}{\sqrt2}\simeq-2.221 .
\end{equation}
\begin{verify}
Finite difference of the numerically solved line: $-2.22144$, against $-\pi/\sqrt2=-2.22144$.
\end{verify}

\subsection{Heterogeneous graphs: the exact rank-one condition}

For $K_{pr}=\rho_r\kappa_p\kappa_r$ (Sec.~\ref{sec:hdmft}), $B_\kappa=\kappa\bar B$ and $q_\kappa=\kappa Q$, so
\begin{equation}
\Delta_\kappa=\frac{a+\mu\kappa\bar B}{\sigma\sqrt{\kappa Q}}=A\kappa^{-1/2}+M\kappa^{1/2},\qquad A=\frac{a}{\sigma\sqrt Q},\quad M=\frac{\mu\bar B}{\sigma\sqrt Q}.
\end{equation}
The closure, Eq.~\eqref{eq:staticclosure}, becomes:
\begin{itemize}
\item $\bar B$ equation: $\bar B=\langle\kappa\cdot\sigma\sqrt{\kappa Q}\,w_1\rangle=\sigma\sqrt Q\langle\kappa^{3/2}w_1\rangle$.
  Multiplying by
  $\mu/(\sigma\sqrt Q)$: $M=\mu\langle\kappa^{3/2}w_1(\Delta_\kappa)\rangle$.
\item $Q$ equation: $Q=\langle\kappa\cdot\sigma^2\kappa Q\,w_2\rangle$, i.e.\ $\sigma^2\langle\kappa^2w_2(\Delta_\kappa)\rangle=1$.
\item normalisation: $\sigma\sqrt Q\,\langle\kappa^{1/2}w_1(\Delta_\kappa)\rangle=1$.
\end{itemize}
The stability matrix $\sigma^2\varphi_pK_{pr}=(\sigma^2w_0(\Delta_p)\kappa_p)(\rho_r\kappa_r)$ is an outer product $u\,v^{\sf T}$; its only
non-zero eigenvalue is $v^{\sf T}u=\sigma^2\langle\kappa^2w_0\rangle$.  At the threshold this equals one, and so does
$\sigma^2\langle\kappa^2w_2\rangle$; subtracting and using $w_2-w_0=\Delta w_1$:
\begin{equation}
\boxed{\sigma_c^{-2}=\langle\kappa^2w_0(\Delta_\kappa)\rangle,\qquad\langle\kappa^2\Delta_\kappa w_1(\Delta_\kappa)\rangle=0,\qquad M=\mu\langle\kappa^{3/2}w_1(\Delta_\kappa)\rangle.}
\label{eq:rankone}
\end{equation}
Three equations for $(A,M,\sigma_c)$.  The normalisation only fixes the scale $Q$ (and with it $a$) and is not
needed for $\sigma_c$.  \controlled

\subsection{Expansion to first order in the degree variance}

Write $\kappa=1+\epsilon$ with $\langle\epsilon\rangle=0$, $\langle\epsilon^2\rangle=V$, and count $\epsilon=O(V^{1/2})$.  We keep
$O(V)$.  On the regular graph ($V=0$), the second relation gives $\Delta=A+M=0$, and the third
$M_0=\mu w_1(0)=\mu\phi_0$, so $A_0=-M_0$.

\paragraph{Step 1: expand $\Delta_\kappa$.}  With $(1+\epsilon)^{-1/2}=1-\tfrac\epsilon2+\tfrac{3\epsilon^2}8$ and
$(1+\epsilon)^{1/2}=1+\tfrac\epsilon2-\tfrac{\epsilon^2}8$,
\begin{equation}
\Delta_\kappa=(A+M)+\frac{M-A}{2}\,\epsilon+\frac{3A-M}{8}\,\epsilon^2 .
\end{equation}
Let $\delta\equiv A+M=O(V)$.  Then $(M-A)/2=M-\delta/2$ and $(3A-M)/8=-M/2+3\delta/8$.  Products $\delta\epsilon$ average
to zero and $\delta\epsilon^2$ is $O(V^2)$, so to the order needed
\begin{equation}
\Delta_\kappa\simeq\delta+M\epsilon-\frac M2\epsilon^2,\qquad \langle\Delta_\kappa\rangle=\delta-\frac{MV}2,\qquad \Var\Delta_\kappa\simeq M^2V.
\end{equation}

\paragraph{Step 2: the second relation.}  Let $G(\Delta)=\Delta w_1(\Delta)$.  Then $G(0)=0$,
$G'=w_1+\Delta w_0$ so $G'(0)=\phi_0$, and $G''=2w_0+\Delta\phi$ so $G''(0)=1$.  To second order,
$G(\Delta_\kappa)\simeq\phi_0(\delta+M\epsilon-\tfrac M2\epsilon^2)+\tfrac12M^2\epsilon^2$.  With $\kappa^2=1+2\epsilon+\epsilon^2$:
\begin{equation}
\langle\kappa^2G\rangle=\phi_0\delta+\phi_0\Big(2MV-\frac{MV}2\Big)+\frac{M^2V}2=\phi_0\delta+\frac32\phi_0MV+\frac12M^2V=0,
\end{equation}
so
\begin{equation}
\phi_0\delta=-\frac32\phi_0MV-\frac12M^2V .
\label{eq:delta}
\end{equation}

\paragraph{Step 3: the Perron root.}  $w_0(\Delta)\simeq\tfrac12+\phi_0\Delta+O(\Delta^3)$ (no quadratic term).  So
\begin{equation}
\langle\kappa^2w_0\rangle=\Big\langle(1+2\epsilon+\epsilon^2)\Big(\tfrac12+\phi_0\big(\delta+M\epsilon-\tfrac M2\epsilon^2\big)\Big)\Big\rangle
=\frac12+\frac V2+\phi_0\delta+\frac32\phi_0MV .
\end{equation}
Insert Eq.~\eqref{eq:delta}: the two terms odd in $\mu$ ($\pm\tfrac32\phi_0MV$) cancel, and
\begin{equation}
\sigma_c^{-2}=\langle\kappa^2w_0\rangle=\frac12\big[1+V(1-M^2)\big],\qquad M^2=M_0^2+O(V)=\frac{\mu^2}{2\pi}.
\end{equation}
Taking the inverse square root,
\begin{equation}
\boxed{\sigma_c\simeq\sqrt2\Big[1-\frac V2\Big(1-\frac{\mu^2}{2\pi}\Big)\Big].}
\label{eq:sigmac}
\end{equation}
\controlled{} to $O(V)$.

\paragraph{Reading.}  The $+V/2$ in the Perron root is the hubs' larger amplitude ($\kappa^2$ weight); it lowers $\sigma_c$.
The $-M^2V/2$ comes only through $\delta$, from $G''(0)=1$: the spread $\Var\Delta_\kappa=M^2V$ of distances from extinction,
acting on the convex function $\Delta w_1$, forces the mean of $\Delta_\kappa$ down to keep $\langle\kappa^2\Delta w_1\rangle=0$, and
since $w_0$ is linear at $0$ this lowers the Perron root.  A mean interaction of either sign therefore raises $\sigma_c$, and
only $\mu^2$ enters at this order.

\begin{verify}
Solving Eq.~\eqref{eq:rankone} exactly for a gamma law of $\kappa$ with mean $1$ and variance $V$
(4000 quantile nodes):
\begin{center}
\begin{tabular}{lcccc}
\toprule
$\mu$ & $V$ & exact & Eq.~\eqref{eq:sigmac} & (exact$-$first order)$/V^2$\\
\midrule
$0$ & $0.01$ & $1.40720$ & $1.40714$ & $0.57$\\
$0$ & $0.04$ & $1.38677$ & $1.38593$ & $0.52$\\
$-1.76$ & $0.01$ & $1.41068$ & $1.41063$ & $0.54$\\
$-1.76$ & $0.04$ & $1.40068$ & $1.39987$ & $0.50$\\
$+1$ & $0.04$ & $1.39126$ & $1.39043$ & $0.52$\\
$-1$ & $0.04$ & $1.39134$ & $1.39043$ & $0.57$\\
\bottomrule
\end{tabular}
\end{center}
The residual scales as $V^2$, as it should.  The small difference between $\mu=\pm1$ at $V=0.04$ is the
odd-in-$\mu$ dependence that appears beyond first order (the SM quotes it at $\alpha=3.5$).
\end{verify}

\subsection{The zero-growth line on heterogeneous graphs}

Zero growth is $a=1$.  Then $A=1/(\sigma\sqrt Q)$, and the normalisation $\sigma\sqrt Q\langle\kappa^{1/2}w_1\rangle=1$ becomes
$\langle\kappa^{1/2}w_1\rangle=A$.  Unknowns $(A,M,\mu)$ at given $(\sigma,V)$:
\begin{equation}
\text{(i)}\ \sigma^2\langle\kappa^2w_2\rangle=1,\qquad \text{(ii)}\ M=\mu\langle\kappa^{3/2}w_1\rangle,\qquad \text{(iii)}\ \langle\kappa^{1/2}w_1\rangle=A .
\end{equation}
At $V=0$: $\Delta=A+M=\Delta_0(\sigma)$, $A=w_1(\Delta_0)$, $M=\mu w_1(\Delta_0)$, which reproduces $\mu^{(0)}=\Delta_0/w_1(\Delta_0)-1$.

\paragraph{General expansion.}  With $\Delta_\kappa=\bar\Delta+d_1\epsilon+d_2\epsilon^2$, $\bar\Delta=A+M$, $d_1=(M-A)/2$,
$d_2=(3A-M)/8$, and $\kappa^n=1+n\epsilon+\tfrac{n(n-1)}2\epsilon^2$, for any smooth $g$:
\begin{equation}
\langle\kappa^ng(\Delta_\kappa)\rangle\simeq g(\bar\Delta)+V\Big[\frac{n(n-1)}2g+n\,d_1g'+d_2g'+\frac{d_1^2}2g''\Big]_{\bar\Delta}.
\label{eq:genexp}
\end{equation}

\paragraph{At $\sigma=\sqrt2$.}  Here $\Delta_0=0$, $A_0=\phi_0$, $M_0=-\phi_0$, $\mu_0=-1$, so $d_1=-\phi_0$ and $d_2=\phi_0/2$.  Let
$\bar\Delta=\delta$ and use $w_2'=2w_1$, $w_2''=2w_0$, $w_1'=w_0$, $w_1''=\phi$.
\begin{itemize}
\item (i) with $n=2$: $w_2(\delta)\simeq\tfrac12+2\phi_0\delta$, and the bracket is
  $w_2+4d_1w_1+2d_2w_1+d_1^2w_0=\tfrac12-4\phi_0^2+\phi_0^2+\tfrac12\phi_0^2=\tfrac12-\tfrac52\phi_0^2$.  So
  $2[\tfrac12+2\phi_0\delta+V(\tfrac12-\tfrac52\phi_0^2)]=1$, i.e.
  \begin{equation}
  \delta=-\frac{V(1-5\phi_0^2)}{4\phi_0}.
  \end{equation}
\item (iii) with $n=\tfrac12$: the bracket is $-\tfrac18w_1+\tfrac12d_1w_0+d_2w_0+\tfrac12d_1^2\phi=-\tfrac{\phi_0}8-\tfrac{\phi_0}4+\tfrac{\phi_0}4+\tfrac{\phi_0^3}2\equiv\beta$, so
  $A=\phi_0+\tfrac\delta2+V\beta$.
\item (ii) with $n=\tfrac32$: the bracket is $\tfrac38w_1+\tfrac32d_1w_0+d_2w_0+\tfrac12d_1^2\phi=\tfrac{3\phi_0}8-\tfrac{3\phi_0}4+\tfrac{\phi_0}4+\tfrac{\phi_0^3}2=\beta$ as
  well.  So $\langle\kappa^{3/2}w_1\rangle=\langle\kappa^{1/2}w_1\rangle=A$ to this order, and (ii) reads $M=\mu A$.
\end{itemize}
With $M=\delta-A=-\phi_0+\tfrac\delta2-V\beta$,
\begin{equation}
\mu=\frac MA=-\frac{\phi_0-\delta/2+V\beta}{\phi_0+\delta/2+V\beta}\simeq-1+\frac{\delta}{\phi_0}
=-1-\frac{V(1-5\phi_0^2)}{4\phi_0^2}=-1+V\Big(\frac54-\frac\pi2\Big),
\end{equation}
using $\phi_0^{-2}=2\pi$.  So $c(\sqrt2)=\tfrac54-\tfrac\pi2\simeq-0.321$: with heterogeneous degrees zero growth needs slightly
stronger competition.  At $\sigma\to0$ all firms survive, $S=a+\mu\kappa\bar B$, and the normalisation with $a=1$ gives
$\mu\bar B=0$ while $\bar B=1/(1-\mu(1+V))\neq0$, so $\mu^{(0)}=0$ for every $V$: $c(0)=0$.  \controlled

\paragraph{Converting the shift to $\sigma$.}  The heterogeneous line is $\mu=\mu^{(0)}_{\rm reg}(\sigma)+cV$.  At fixed $\mu$ its $\sigma$
moves by $-cV/(\dd\mu^{(0)}/\dd\sigma)=-(-0.321V)/(-2.221)=-0.144V$.  The shift of $\sigma_c$ at $\mu=-1$ from Eq.~\eqref{eq:sigmac} is
$-\tfrac{\sqrt2}2(1-\tfrac1{2\pi})V=-0.595V$.  Ratio $0.24$: ``about a quarter''.

\begin{verify}
Numerical solution of (i)--(iii) with a gamma law of $\kappa$: the shift in $\mu$ at $\sigma=\sqrt2$ divided by $V$ is $-0.319$
($V=0.01$) and $-0.318$ ($V=0.02$), against $-0.321$.  At other $\sigma$, $c=-0.056$, $-0.153$, $-0.235$ for $\sigma=0.7$,
$1.0$, $1.2$: monotone from $0$ to $-0.32$, as stated.
\end{verify}

%======================================================================
\section{One node with demographic noise at constant field}
\label{sec:constant}

\subsection{Scale function, speed measure and boundaries}

At constant $\Phi$, Eq.~\eqref{eq:onefirm} is the one-dimensional diffusion
$\dd S=b(S)\,\dd t+\sqrt{a(S)}\,\dd W$ with $b(y)=y(\Phi-y)$ and $a(y)=Dy$, generator
$\mathcal L=\tfrac12Dy\,\partial_y^2+y(\Phi-y)\partial_y$.  The scale density and speed density
\cite{karlintaylor1981} are
\begin{equation}
s'(y)=\exp\Big[-\int_0^y\frac{2b(u)}{a(u)}\dd u\Big]=\exp\Big[-\frac2D\Big(\Phi y-\frac{y^2}2\Big)\Big],\qquad
m(y)=\frac{2}{a(y)s'(y)}=\frac{2}{Dy\,s'(y)},
\end{equation}
since $2b/a=2(\Phi-u)/D$.  Let $s(y)=\int_0^ys'(u)\,\dd u$.  Properties:
\begin{itemize}
\item Near $0$, $s(y)\simeq y$ and $m(y)\simeq2/(Dy)$; $\int_0s(y)m(y)\,\dd y<\infty$, so $0$ is reached in finite time
  (an exit boundary in Feller's classification).
\item $s'(y)\sim e^{y^2/D}$ for large $y$, so $s(\infty)=\infty$: $\infty$ is never reached, and the process started anywhere
  hits $0$ with probability one.  Every life ends, so at a node with constant field the lives form a proper
  renewal sequence.
\item $\int^\infty m(z)\,\dd z<\infty$ (Gaussian decay).
\end{itemize}
\exact

\subsection{Establishment probability}

For $0<y_1<y_0<y_2$, $\Pr(\text{hit }y_2\text{ before }y_1\mid y_0)=[s(y_0)-s(y_1)]/[s(y_2)-s(y_1)]$ (the scale function is
the harmonic function of $\mathcal L$).  With $y_1=0$, $y_0=D$, $y_2=\Phi$:
\begin{equation}
P_{\rm est}=\frac{\int_0^Ds'(y)\,\dd y}{\int_0^\Phi s'(y)\,\dd y}\qquad(\Phi>D).
\end{equation}
\exact{}  \emph{Feller comparison:} dropping the logistic term, $s'(y)=e^{-2\Phi y/D}$, and the probability of
escaping to large sizes from $S_0=D$ is $1-e^{-2\Phi S_0/D}=1-e^{-2\Phi}$, independent of $D$.

\subsection{Mean exit time}

$T(x)$ solves $\mathcal LT=-1$ with $T(0)=0$ and no flux at infinity.  Since $s''/s'=-2b/a$,
\begin{equation}
\frac{\dd}{\dd y}\Big(\frac{T'}{s'}\Big)=\frac{T''+(2b/a)T'}{s'}=\frac{2\,\mathcal LT}{a\,s'}=-\frac{2}{a\,s'}=-m(y).
\end{equation}
Integrating from $y$ to $\infty$ with $T'/s'\to0$ at infinity, $T'(y)=s'(y)\int_y^\infty m(z)\,\dd z$, and
\begin{equation}
T(x)=\int_0^xs'(y)\int_y^\infty m(z)\,\dd z\,\dd y=2\int_0^xs'(y)\int_y^\infty\frac{\dd z}{Dz\,s'(z)}\,\dd y .
\label{eq:T}
\end{equation}
\exact

\subsection{Green's function and the renewal law}

Exchanging the order of integration in Eq.~\eqref{eq:T} (Fubini; all integrands positive),
\begin{equation}
T(x_0)=\int_0^\infty m(z)\int_0^{\min(x_0,z)}s'(y)\,\dd y\,\dd z=\int_0^\infty\underbrace{s(\min(x_0,z))\,m(z)}_{G(x_0,z)}\dd z .
\end{equation}
$G(x_0,z)$ is the Green's function: the expected time spent near $z$, per unit $z$, before absorption, starting at $x_0$.

\paragraph{Renewal--reward.}  With replacement at $D$, the successive lives at a node are i.i.d.\ copies of the
process started at $D$.  For any set $\mathcal A$, the long-run fraction of time spent in $\mathcal A$ equals
$\E[\text{time in }\mathcal A\text{ during one life}]/\E[\text{life length}]$.  So the stationary density of sizes at the node is
\begin{equation}
\boxed{p(x\mid h)=\frac{G(D,x)}{T(D)}=\frac{2\,s(\min(x,D))}{D\,x\,s'(x)\,T(D)},\qquad T(D)=\int_0^\infty G(D,x)\,\dd x,}
\label{eq:renewal}
\end{equation}
and by the elementary renewal theorem the turnover rate is $1/T(D)$ per node per unit time.  \exact

\paragraph{Shape.}  For $x>D$, $s(\min(x,D))=s(D)$ is constant and
$1/(xs'(x))=x^{-1}\exp[(2/D)(hx-x^2/2)]=x^{-1}e^{h^2/D}e^{-(x-h)^2/D}$: a Gaussian factor of variance $D/2$ around
$h$ times $x^{-1}$.  For $x<D$ and $x\ll\sqrt D$, $s(x)\simeq x$ and $s'(x)\simeq1$, so $p\simeq2/(DT(D))$: flat.  \exact

\begin{verify}
At $h=0.5$, $D=0.1$, $\int G(D,x)\,\dd x=22.8305$ and the double integral \eqref{eq:T} give the same value to $10^{-12}$.
\end{verify}

\subsection{Asymptotics of the exit time}

\paragraph{(a) Negative field, $-\Phi\gg\sqrt D$: $T(D)$ does not depend on $D$.}  Rescale $y=D\tilde y$, $z=D\tilde z$ in
Eq.~\eqref{eq:T} with $x=D$.  Then $s'(D\tilde y)=e^{2|\Phi|\tilde y+D\tilde y^2}$, $\dd y=D\,\dd\tilde y$, and
$\dd z/(Dz\,s'(z))=\dd\tilde z/(D\tilde z\,s'(D\tilde z))$.  So
\begin{equation}
\begin{aligned}
T(D)&=2\int_0^1e^{2|\Phi|\tilde y+D\tilde y^2}\int_{\tilde y}^\infty\frac{e^{-2|\Phi|\tilde z-D\tilde z^2}}{\tilde z}\dd\tilde z\,\dd\tilde y\\
&\xrightarrow{D\to0}\;2\int_0^1e^{2|\Phi|\tilde y}E_1(2|\Phi|\tilde y)\,\dd\tilde y
=\frac1{|\Phi|}\int_0^{2|\Phi|}e^tE_1(t)\,\dd t ,
\end{aligned}
\end{equation}
with $E_1$ the exponential integral.  The $D\tilde z^2$ terms matter only for $\tilde z\gtrsim D^{-1/2}$, which the factor
$e^{-2|\Phi|\tilde z}$ suppresses when $|\Phi|\gg\sqrt D$.  The limit is the Feller (linear birth--death) value.  Since $e^tE_1(t)\sim1/t$,
$T\sim\ln(2|\Phi|)/|\Phi|$ for large $|\Phi|$: of order $1/|\Phi|$ up to a logarithm.  \controlled

\begin{verify}
At $\Phi=-0.5$ the limit is $2.3471$, which coincides with $\int_0^\infty[1-\exp(-2|\Phi|/(e^{|\Phi|t}-1))]\dd t=2.3471$ from the
Feller extinction law; the exact quadrature gives $2.323$, $2.156$, $1.924$ at $D=0.01$, $0.1$, $0.3$ (SM: $2.32$, $2.16$, $1.92$).
\end{verify}

\paragraph{(b) Zero field: $T(D)=\ln(1/D)+2-\gamma_E+O(D)$.}  At $\Phi=0$, $s'(y)=e^{y^2/D}$ and
$m(z)=(2/Dz)e^{-z^2/D}$.  For $y\le D$, $y^2/D\le D$, so $s'(y)=1+O(D)$.  The inner integral, with $t=z^2/D$
($\dd z/z=\dd t/2t$), is
\begin{equation}
\int_y^\infty\frac{2}{Dz}e^{-z^2/D}\dd z=\frac1DE_1\Big(\frac{y^2}{D}\Big)=\frac1D\Big[-\gamma_E-\ln\frac{y^2}{D}+O\Big(\frac{y^2}D\Big)\Big].
\end{equation}
Then
\begin{equation}
T(D)=\frac1D\int_0^D\Big[-\gamma_E+\ln D-2\ln y\Big]\dd y=-\gamma_E+\ln D-2(\ln D-1)=\ln\frac1D+2-\gamma_E .
\end{equation}
The article says ``$T$ grows as $\ln(1/D)$''; this is the exact asymptotic with its constant, and it also fixes the
coefficient of the logarithm to $1$.  The heuristic: the critical Feller survival $\Pr(T>t)\simeq2S_0/(Dt)=2/t$ is cut off by the
logistic term at $t\sim D^{-1/2}$, giving $\int_1^{D^{-1/2}}2\,\dd t/t=\ln(1/D)$.  \controlled

\begin{verify}
$T(D)-\ln(1/D)=1.442$, $1.425$, $1.423$ at $D=10^{-2},10^{-3},10^{-4}$, against $2-\gamma_E=1.4228$.  At $D=0.01$, $0.1$, $0.3$ the
quadrature gives $6.047$, $3.842$, $2.881$ (SM: $6.05$, $3.84$, $2.88$).
\end{verify}

\paragraph{(c) Positive field: $T\sim e^{\Phi^2/D}$.}  For a firm on its plateau, $s'(y)$ falls from $1$ at $y=0$ over a scale $D/\Phi$,
so $\int_0^{S_0}s'\simeq D/(2\Phi)$; $\int_y^\infty m$ is dominated by the peak of $m$ at $z\simeq\Phi$, of height
$(2/D\Phi)e^{\Phi^2/D}$ and width $\sqrt{\pi D}$ (Laplace).  So
\begin{equation}
T\simeq\frac{\sqrt{\pi D}}{\Phi^2}\,e^{\Phi^2/D}\big[1+O(D/\Phi^2)\big].
\end{equation}
\controlled{} in $D/\Phi^2$.  An established firm exits through dips of its field rather than through its own noise.
\begin{verify}
Exact against Laplace at $D=0.1$: $13.18$ vs $15.3$ ($\Phi=0.3$), $38.0$ vs $27.3$ ($\Phi=0.5$), $200$ vs $154$ ($\Phi=0.7$).  The exponential
order is right, the prefactor is off by 15--40\% at $\Phi^2/D=1$--$5$; the article only claims the exponential order.
\end{verify}

\subsection{The plateau as an Ornstein--Uhlenbeck process}

Linearise $S=\Phi+\varepsilon$: $\dd\varepsilon=-\Phi\varepsilon\,\dd t+\sqrt{D\Phi}\,\dd W$, an OU process with rate $\Phi$ and stationary
variance $D\Phi/(2\Phi)=D/2$.  In logs, $\ln S-\ln\Phi\simeq\varepsilon/\Phi$ has rate $\Phi$ and variance $D/(2\Phi^2)$.  For an OU
process of rate $k$ and stationary variance $v$, the increment over $\Delta t$ has variance $2v(1-e^{-k\Delta t})$, so
\begin{equation}
\Var g=\frac{D}{\Phi^2}\big(1-e^{-\Phi\Delta t}\big).
\label{eq:ouvar}
\end{equation}
\controlled{} in $D/\Phi^2$.  (Including the $-D/2S$ It\^o drift of Eq.~\eqref{eq:log} shifts the rate to $\Phi-D/2\Phi$, a
correction of the same order.)

\begin{verify}
Simulation of one node at $D=0.1$, absorbing at $0$, growth rates within lives only, $\Delta t=0.5$:
\begin{center}
\begin{tabular}{lccc}
\toprule
$\Phi$ & simulated $\Var g$ & Eq.~\eqref{eq:ouvar} & ratio\\
\midrule
$0.7$ & $0.104$ & $0.060$ & $1.72$\\
$1.0$ & $0.047$ & $0.039$ & $1.20$\\
$1.5$ & $0.025$ & $0.024$ & $1.08$\\
$2.0$ & $0.0165$ & $0.0158$ & $1.04$\\
$4.0$ & $0.0055$ & $0.0054$ & $1.02$\\
\bottomrule
\end{tabular}
\end{center}
At $D=0.02$ the ratio is $1.02$ already at $\Phi=1$.  The excess behaves as $1+c\,D/\Phi^2$ with $c\simeq2$.
\end{verify}
\begin{flag}
SM Sec.~S3 says the OU formula is ``accurate for $\Phi\gtrsim1$''.  The variance error is about $2D/\Phi^2$: at $D=0.1$ the
formula is $20\%$ low in variance ($10\%$ in standard deviation) at $\Phi=1$, and $5\%$ in standard deviation needs
$D/\Phi^2\lesssim0.05$ ($\Phi\gtrsim1.4$).  Suggest ``accurate to $10\%$ in standard deviation for $\Phi\gtrsim1$ at $D=0.1$''.  The same
nonlinearity is what makes the volatility law worst in the lowest size decile (Sec.~\ref{sec:law}).
\end{flag}

%======================================================================
\section{The static closure with noise}
\label{sec:noisyclosure}

When the field is static, each node follows the renewal law \eqref{eq:renewal} at its own $h=a+\mu B_p+\sigma\sqrt{q_p}\,z$.
Let $\langle S\rangle_h=\int xp(x\mid h)\,\dd x$.  The static closure is Eq.~\eqref{eq:closure} with time averages:
\begin{equation}
B_p=\sum_rK_{pr}\E\langle S\rangle_{h_r},\qquad q_p=\sum_rK_{pr}\E\langle S\rangle_{h_r}^2,\qquad \sum_p\rho_p\E\langle S\rangle_{h_p}=1,
\end{equation}
with $\E$ the Gaussian average over $z$.  Using $\langle S\rangle^2_h$ in $q_p$ counts only the static part of the neighbours'
sizes in the field covariance, i.e.\ the temporal fluctuations of $S$ produced by the noise are taken as not
feeding the field.  (Using $\langle S^2\rangle_h$ instead, ``as if quenched'', moves the lines by less than $0.015$; \numer.)

\paragraph{Stability.}  Repeat the response argument of Sec.~\ref{sec:static} with $\delta\langle S\rangle_h=\chi(h)\,\delta h$,
$\chi=\partial\langle S\rangle_h/\partial h$.  The variance recursion becomes $Y_p=\E[\chi_p^2](\epsilon+\sigma^2\sum_rK_{pr}Y_r)$, so the
threshold is the Perron root of $\sigma^2{\rm diag}(\E\chi^2)K$ reaching one.  Without noise $\langle S\rangle_h=[h]_+$, $\chi=\theta(h)$,
$\chi^2=\chi$ and $\E\chi^2=\varphi$: the deterministic criterion is recovered.  \controlled{} (static field).

\paragraph{The floor version.}  With $\lambda(1-S)$ and no noise, the stationary size solves $S(h-S)+\lambda(1-S)=0$, i.e.\
$S^2-(h-\lambda)S-\lambda=0$, with positive root $S=\tfrac12[(h-\lambda)+\sqrt{(h-\lambda)^2+4\lambda}]$.  Differentiating the
equation, $\chi(h-2S-\lambda)+S=0$; from the equation itself $h=S+\lambda-\lambda/S$, so $2S+\lambda-h=S+\lambda/S$ and
\begin{equation}
\chi=\frac{S}{S+\lambda/S}=\frac{S^2}{S^2+\lambda}.
\end{equation}
\exact{}  $\chi$ switches from $0$ to $1$ over $h\sim\sqrt\lambda$, hence the dense grid in the solver.

\paragraph{Turnover.}  $\sum_p\rho_p\E[1/T(D;h_p)]$ per node per unit time.  Nodes with $h<0$ dominate it, and by
Sec.~\ref{sec:constant}(a) their newborns' lives are nearly $D$-independent: turnover is nearly independent of the noise.

\paragraph{What is not in the closure.}  A static closure cannot produce the collective amplification of slow field
fluctuations near $\sigma_c$ (it underestimates the temporal share of the field variance by factors $2$--$5$ below $\sigma_c$,
\numer).  Above $\sigma_c$ it is a continuation, not a controlled solution; this is why the zero-growth line there is drawn
dashed.

%======================================================================
\section{The deterministic limit: life cycles}
\label{sec:lifecycle}

The floor model has $\dd\ln S_i/\dd t=\Phi_i-S_i+\lambda(S_i^{-1}-1)$, and a life (``spell'') is a maximal stretch above
$S_{\rm thr}=8\lambda$.

\subsection{Size is mean fitness}

Over a window where $\ln S_i$ returns to its start, $0=\bar\Phi_i-\Sb_i+\lambda(\overline{S_i^{-1}}-1)$, so
$\Sb_i=\bar\Phi_i+\lambda(\overline{S_i^{-1}}-1)\simeq\bar\Phi_i$.  \exact{} (the first equality).

\subsection{The fitness law}

From Eq.~\eqref{eq:fieldlaw}, $\Phi(t)=r_\kappa+\sigma\sqrt\kappa\,\xi(t)$ with $r_\kappa=1-\fb+\mu\kappa\bar B$ and
$\E[\xi(t)\xi(t')]=Q(t-t')$ in a stationary state.  Write $Q_T=\lim_{L\to\infty}Q(L)$: the part of the
correlation that does not decay.  It is non-zero because $\xi$ is built from the neighbours' sizes through
\emph{quenched} couplings; the neighbours' time-averaged sizes $\Sb_j>0$ produce a static random field, and
$Q_T=\sum_jw_j\Sb_j^2$ with degree weights $w_j$.  Decompose $\xi=\xi_{\rm s}+\xi_{\rm t}$ into a static part of variance
$Q_T$ and a temporal part of variance $Q_0-Q_T$ with correlation $\rho(L)=(Q(L)-Q_T)/(Q_0-Q_T)$.  Over a long window the
time mean of $\Phi$ is $r_\kappa+\sigma\sqrt\kappa\,\xi_{\rm s}$, so
\begin{equation}
\bar\Phi\mid\kappa\sim\mathcal N(r_\kappa,a_\kappa^2),\quad a_\kappa=\sigma\sqrt{\kappa Q_T},\qquad f_\kappa=\sigma\sqrt{\kappa(Q_0-Q_T)},
\end{equation}
with $f_\kappa$ the temporal amplitude.  \controlled{} (for a finite window $Q_T$ is the window-static part, which
includes slow components; that is the meaning of the subscript $T$).

\subsection{The exact one-firm solution at $\lambda=0$}

With $u=1/S$, $\dot u=-\dot S/S^2=-(\Phi-S)/S=-\Phi u+1$, a linear equation.  With $\mathcal A(t)=\int^t\Phi$ and the
initial condition sent to $-\infty$,
\begin{equation}
\frac1{S(t)}=\int_{-\infty}^te^{-[\mathcal A(t)-\mathcal A(t')]}\,\dd t' .
\end{equation}
\exact{}  (With a finite initial time $t_0$, add $u(t_0)e^{-[\mathcal A(t)-\mathcal A(t_0)]}$.)

\paragraph{Near a downcrossing.}  Let $\Phi$ cross zero downward at $t_0$ with slope $-s$, $\Phi\simeq-s(t-t_0)$.  Then
$\mathcal A(t_0)-\mathcal A(t')=s(t'-t_0)^2/2\ge0$: $t_0$ is a local maximum of $\mathcal A$.
\begin{itemize}
\item At $t=t_0$: $1/S(t_0)=\int_{-\infty}^{t_0}e^{-s(t'-t_0)^2/2}\dd t'=\tfrac12\sqrt{2\pi/s}$, so $S(t_0)=\sqrt{2s/\pi}$.
\item For $t-t_0\gg s^{-1/2}$: $1/S(t)=e^{\mathcal A(t_0)-\mathcal A(t)}\int_{-\infty}^te^{-[\mathcal A(t_0)-\mathcal A(t')]}\dd t'\simeq e^{\mathcal A(t_0)-\mathcal A(t)}\sqrt{2\pi/s}$
  (the Gaussian integral is now complete), so
  \begin{equation}
  S(t)\simeq\sqrt{\frac{s}{2\pi}}\;e^{-[\mathcal A(t_0)-\mathcal A(t)]}.
  \end{equation}
  For intermediate times the incomplete Gaussian integral gives the inverse Mills ratio.
\end{itemize}
\controlled{} (requires $\Phi$ linear over a time $s^{-1/2}$ around $t_0$).

\paragraph{Kill rule.}  $\mathcal A(t_0)-\mathcal A(t)=\int_{t_0}^t(-\Phi)$ is the area of the negative excursion so far.  $S$ drops below
$S_{\rm thr}$ once
\begin{equation}
\int_{t_0}^t(-\Phi)\,\dd t'>\ell_c(s)=\ln\frac{\sqrt{s/2\pi}}{S_{\rm thr}} ;
\end{equation}
if the excursion ends first, $S$ recovers.  A spell ends at the first negative excursion whose area exceeds $\ell_c$.
\heur{} in that it neglects the floor and dips that start before $S$ has recovered from the previous one
(96--98\% correct, \numer).

\subsection{Rice rate and the lifetime law}

Write $\Phi=f(u+\eta)$ with $\eta$ stationary Gaussian of unit variance and correlation $\rho$, so $\Phi=0$ is $\eta=-u$.  Rice's
formula gives the rate of downcrossings of a level $c$ as $\nu=\frac1{2\pi}\sqrt{-\rho''(0)}\,e^{-c^2/2}$.  For
$\rho(L)=(1+L/\tau)e^{-L/\tau}$, expanding with $x=L/\tau$: $(1+x)(1-x+\tfrac{x^2}2-\tfrac{x^3}6)=1-\tfrac{x^2}2+\tfrac{x^3}3+\dots$, so
$\rho''(0)=-1/\tau^2$ and
\begin{equation}
\nu(u)=\frac{e^{-u^2/2}}{2\pi\tau}.
\end{equation}
\exact{} (given the Gaussian process).  A fraction $\PhiN(u)$ of time has $\Phi>0$.  If a firm is listed when its field is
positive and each downcrossing kills with probability $P_{\rm kill}$, the mean spell is the listed time per kill:
\begin{equation}
\bar L(u)=\frac{\PhiN(u)}{\nu(u)P_{\rm kill}(u)} .
\end{equation}
\heur{} (listed time $\approx$ time with $\Phi>0$; entry and exit delays neglected, which is why it runs 10--20\% low for $u\le0.5$).

\subsection{Large $u$: the Slepian parabola}

Condition on a downcrossing of $-u$ at $t=0$ with slope $\eta'(0)=-\dot s$.  For a stationary Gaussian process,
$\E[\eta''(0)\mid\eta(0)=c]=\rho''(0)\,c$, so with $c=-u$ the conditional curvature is $u/\tau^2$, and near the crossing
\begin{equation}
\eta(t)+u\simeq-\dot st+\frac{u}{2\tau^2}t^2 .
\end{equation}
The excursion below $-u$ ends at $t^*=2\dot s\tau^2/u$, and its area in units of $\Phi$ is
\begin{equation}
f\int_0^{t^*}\Big(\dot st-\frac{u t^2}{2\tau^2}\Big)\dd t=f\,t^{*2}\Big(\frac{\dot s}2-\frac{ut^*}{6\tau^2}\Big)=f\,t^{*2}\frac{\dot s}{6}=\frac23\,\frac{f\dot s^3\tau^4}{u^2}.
\end{equation}
The kill condition $\frac23f\dot s^3\tau^4/u^2>\ell_c$ is $\dot s>s_*=(3\ell_cu^2/2f\tau^4)^{1/3}$.  Crossing slopes follow the Rayleigh law
$p(\dot s)\propto\dot s\,e^{-\dot s^2\tau^2/2}$ (variance of $\eta'$ is $-\rho''(0)=\tau^{-2}$), so $\Pr(\dot s>s_*)=e^{-s_*^2\tau^2/2}$ and
\begin{equation}
-\ln P_{\rm kill}\simeq\frac{s_*^2\tau^2}{2}=\frac12\Big(\frac{3\ell_c}{2f\tau}\Big)^{2/3}u^{4/3}.
\end{equation}
\controlled{} for $u\gg1$, treating $\ell_c$ (which depends on the slope only through a logarithm) as constant.
\begin{verify}
The parabola area: numerical quadrature equals $\tfrac23f\dot s^3\tau^4/u^2$ to machine precision.
\end{verify}

\paragraph{Consequences.}  $\ln\bar L=\tfrac{u^2}2+\tfrac12(3\ell_c/2f\tau)^{2/3}u^{4/3}+\ln(2\pi\tau)+O(1)$, dominated by $u^2/2$.  With
$\Sb\simeq uf$, $\ln\bar L\simeq\Sb^2/2f^2$: lifetime grows faster than any power of size.  A window $T$ resolves lifetimes up to
$\bar L(u_\times)\simeq T$, i.e.\ $u_\times\simeq\sqrt{2\ln T}$.

\subsection{The lifetime identity}
\label{sec:lifetimeidentity}

\begin{proposition}
Let a spell start and end at $S_{\rm thr}$ with growth rates $g_1,\dots,g_n$ and peak $S_{\max}$.  Then $\bar g=0$ and
$v=\sqrt{\pi/2}\,(2\ell+V_p)/n$ with $\ell=\ln(S_{\max}/S_{\rm thr})$ and $V_p\ge0$.  More generally, for entry size $S_a$ and exit
size $S_b$, $v\ge\sqrt{\pi/2}\;2\ln[S_{\max}/\max(S_a,S_b)]/n$.
\end{proposition}
\begin{proof}
$\sum g_t=\ln(S_b/S_a)$, which is $0$ for $S_a=S_b$.  Split the steps at the peak into $n_1$ before and $n_2$ after.  Let
$L_a=\ln(S_{\max}/S_a)\ge0$ and $L_b=\ln(S_{\max}/S_b)\ge0$, so $n\bar g=L_a-L_b$.  By the triangle inequality on each side,
\begin{equation}
\sum_t|g_t-\bar g|\ge|L_a-n_1\bar g|+|L_b+n_2\bar g| .
\end{equation}
Suppose $L_a\ge L_b$ (so $\max(S_a,S_b)=S_b$ and $\bar g\ge0$).  Then $n_1\bar g\le n\bar g=L_a-L_b$, so $L_a-n_1\bar g\ge L_b$, and
$L_b+n_2\bar g\ge L_b$.  So the sum is at least $2L_b$.  The other case is symmetric.  For $S_a=S_b=S_{\rm thr}$, $\bar g=0$ and the
sum is $2\ell$ plus the excess variation $V_p$.
\end{proof}
\exact{}  The bound holds with the MAD taken around the mean, as the estimator does.

\subsection{The cumulant criterion for multiscaling}
\label{sec:cumulant}

\begin{proposition}
Let $v=A(\Sb)X$ and let $c_n(\Sb)$ be the cumulants of $\ln X$ at fixed size.  Then
\begin{equation}
\zeta_q-q\zeta_1=-\sum_{n\ge2}\frac{q^n-q}{n!}\frac{\dd c_n}{\dd\ln\Sb}=-\frac{q(q-1)}2\frac{\dd\Var(\ln X\mid\Sb)}{\dd\ln\Sb}-\frac{q(q^2-1)}{6}\frac{\dd c_3}{\dd\ln\Sb}-\dots
\end{equation}
\end{proposition}
\begin{proof}
$\ln\E[v^q\mid\Sb]=q\ln A+\ln\E[e^{q\ln X}]=q\ln A+\sum_{n\ge1}c_nq^n/n!$.  Differentiate in $\ln\Sb$ and subtract $q$ times the $q=1$
case; the $n=1$ terms cancel.
\end{proof}
\exact{} (when the cumulant series converges).  To second order the curve bends below $q$ iff the conditional variance of
$\ln v$ grows with size, and lies above $q$ if it shrinks.  This is the general form of every bend statement in this note: the
multiscaling identity of Sec.~\ref{sec:identity} is a special case that identifies \emph{which} dispersion grows.

For a spell, $X\propto\ell/n$.  Large firms exit on rare deep dips, nearly a Poisson event, so $n$ is close to exponential and
$\Var\ln n\to\pi^2/6$ (the logarithm of an exponential variable is Gumbel-distributed with variance $\pi^2/6$).

%======================================================================
\section{The volatility law}
\label{sec:law}

\subsection{Decomposition}

Start from Eq.~\eqref{eq:log}.  Let $y_F$ solve the noise-free equation $\dot y_F=\Phi-e^{y_F}$ along the same field and from
the same initial size, and write $\ln S=y_F+\delta$.  Subtracting,
\begin{equation}
\dd\delta=-e^{y_F}\big(e^{\delta}-1\big)\dd t-\frac{D}{2S}\,\dd t+\sqrt{\frac DS}\,\dd W .
\end{equation}
\begin{itemize}
\item Linearise in $\delta$: $-e^{y_F}(e^\delta-1)\simeq-S_F\delta$ with $S_F=e^{y_F}\simeq S$.
\item Freeze the coefficients over the relaxation time $1/S$ (adiabatic): $\delta$ is an OU process with rate $k=S$ and noise
  intensity $D/S$, stationary variance $v_\delta=(D/S)/(2S)=D/(2S^2)$.
\item The It\^o drift $-D/2S$ shifts the mean of $\delta$, not its fluctuations, to this order.
\end{itemize}
The increment over $\Delta t$ has variance $2v_\delta(1-e^{-k\Delta t})$:
\begin{equation}
\Var(\Delta\delta)=\frac{D}{S^2}\big(1-e^{-S\Delta t}\big).
\end{equation}
The growth rate is $g=\Delta y_F+\Delta\delta$.  Taking the two parts as uncorrelated (the noise $W$ is independent of the field,
and $y_F$ is a functional of the field only) and averaging the local variance over the life:
\begin{equation}
\boxed{\Var g\simeq V_F+B,\qquad B=\Big\langle\frac D{S^2}\big(1-e^{-S\Delta t}\big)\Big\rangle_{\rm life},}
\end{equation}
where $V_F$ is the growth variance of $y_F$ over the life.  Limits: $S\Delta t\ll1$ gives $B\to D\Delta t\langle S^{-1}\rangle$; $S\Delta t\gg1$ gives
$B\to D\langle S^{-2}\rangle$.  \heur.

\paragraph{Conditions.}  (1) $\delta$ small: $v_\delta=D/(2S^2)\ll1$, i.e.\ $S\gg\sqrt{D/2}\simeq0.22$ at $D=0.1$ (on the plateau, Sec.~\ref{sec:constant}
finds a variance error of about $2D/S^2$, i.e.\ $20\%$ at $S=1$; along a life the relevant $S$ is the instantaneous one, and
the error is partly averaged out over a life that also visits larger sizes).  (2) $S$ varies slowly on the time $1/S$.  (3) The cross term $2\Cov(\Delta y_F,\Delta\delta)$ is small: $\delta$
relaxes towards $y_F$, so the deviation and the path are not exactly independent once the linearisation fails.  All three are
worst for small firms, consistent with the 9\% error in the lowest decile ($\Sb=0.26$) against 5\% elsewhere.

\subsection{Why the filter matters}

If a firm followed its field exactly it would sit on $S=\Phi$, with $\Var g=\Var\Delta\ln\Phi$ plus unfiltered noise
$D\Delta t\langle S^{-1}\rangle$.  The noise part changes little (at $\Delta t=0.5$, $\sqrt{D\Delta t\langle S^{-1}\rangle}$ is within 6\% of $\sqrt B$),
but $\Var\Delta\ln\Phi$ is much larger than $V_F$ for small firms: $\ln\Phi$ moves violently when $\Phi$ is small, while $y_F$
low-pass filters the field at rate $S$.  (Numbers \numer: $\sqrt{\Var\Delta\ln\Phi}$ from $0.92$ to $0.44$, $\sqrt{V_F}$ flat at $0.32$--$0.37$.)

\subsection{The crossover}

With $\sqrt{V_F}$ nearly flat and $B\propto\Sb^{-b}$, $v^2\simeq v_F^2+B(\Sb)$ has local slope
$-\tfrac12\dd\ln v^2/\dd\ln\Sb=\tfrac b2\,w$ with $w=B/(v_F^2+B)$: $b/2$ where noise dominates, $0$ where the floor does.  From SM
Table~S4, $\sqrt B$ falls from $0.724$ to $0.260$ while $\Sb$ goes from $0.26$ to $3.53$: $b=2\ln(0.724/0.260)/\ln(3.53/0.26)=0.79$,
within the quoted $0.66$--$0.87$.  The two parts cross between $\Sb=1.43$ ($0.366$ vs $0.381$) and $1.73$ ($0.365$ vs $0.350$), i.e.\
near $1.6$.  \exact{} (arithmetic on the table).

%======================================================================
\section{The multiscaling mechanism}
\label{sec:identity}

\subsection{The identity}

Assume $v_i^2=(u_i+B(\Sb_i))\,\epsilon_i^2$, with $\epsilon_i$ independent of $u_i$ and of $\Sb_i$ (estimation noise), and $B$ a function
of size alone.  Let $p(u\mid\Sb)$ be the law of the floor at fixed size and
\begin{equation}
E_q(\Sb)=\int p(u\mid\Sb)\,(u+B(\Sb))^{q/2}\,\dd u,\qquad \E[v^q\mid\Sb]=E_q(\Sb)\,\E[\epsilon^q].
\end{equation}
Since $\E[\epsilon^q]$ does not depend on size, $\zeta_q=-\dd\ln E_q/\dd\ln\Sb$.  Differentiate under the integral; there are two
contributions, from $B(\Sb)$ inside the power and from $p$:
\begin{equation}
\frac{\dd E_q}{\dd\ln\Sb}=\int p\,\frac q2(u+B)^{q/2-1}\frac{\dd B}{\dd\ln\Sb}\,\dd u+\int\frac{\partial p}{\partial\ln\Sb}(u+B)^{q/2}\dd u .
\end{equation}
With $\dd B/\dd\ln\Sb=-bB$ and $(u+B)^{q/2-1}B=(u+B)^{q/2}w$, $w=B/(u+B)$, the first term is $-\tfrac q2bE_q\langle w\rangle_q$, where
\begin{equation}
\langle w\rangle_q=\frac{\E[(u+B)^{q/2}w\mid\Sb]}{\E[(u+B)^{q/2}\mid\Sb]} .
\end{equation}
Dividing by $-E_q$:
\begin{equation}
\boxed{\zeta_q=\frac q2\,b\,\langle w\rangle_q-c_q,\qquad c_q=\frac1{E_q}\int\frac{\partial p(u\mid\Sb)}{\partial\ln\Sb}(u+B)^{q/2}\,\dd u .}
\end{equation}
\exact{} under the stated assumptions.  $c_q=0$ when the law of the floor does not depend on size.

\paragraph{Floor whose law shifts with size.}  If $\ln u=\delta(\Sb)+\xi$ with the law of $\xi$ fixed, then the size dependence of $p$
is entirely through $\delta$, and $\partial_{\ln\Sb}$ at fixed $B$ acts on $(e^{\delta+\xi}+B)^{q/2}$ as
$\tfrac q2(u+B)^{q/2-1}u\,\delta'$.  So
\begin{equation}
c_q=\frac q2\big(1-\langle w\rangle_q\big)\delta',\qquad
\zeta_q=\frac q2\Big[b\langle w\rangle_q-\big(1-\langle w\rangle_q\big)\delta'\Big],\qquad \delta'=\frac{\dd\delta}{\dd\ln\Sb}.
\end{equation}
A rising floor ($\delta'>0$) flattens the decline.  It also adds to the bend: write $\zeta_q=\tfrac q2[(b+\delta')x_q-\delta']$ with
$x_q=\langle w\rangle_q$.  For $x_q<x_1$ and $\zeta_q,\zeta_1>0$,
\begin{equation}
\frac{(b+\delta')x_q-\delta'}{(b+\delta')x_1-\delta'}<\frac{x_q}{x_1}\iff -x_1\delta'<-x_q\delta'\iff x_q<x_1 ,
\end{equation}
so the ratio is below the $c_q=0$ value $qx_q/x_1$.  \exact

\paragraph{Homogeneous floor.}  If $u$ is the same for all firms of a given size, however it depends on size, then
$E_q=Y(\Sb)^{q/2}$ with $Y=u+B$ and $\zeta_q=-\tfrac q2\dd\ln Y/\dd\ln\Sb=q\zeta_1$ exactly, whatever $c_q$.  \exact

\subsection{Monotonicity of $\langle w\rangle_q$ and the bound $\zeta_q/\zeta_1\le q$}

With $Y=u+B$ at fixed size, $\langle w\rangle_q=\E[Y^{q/2}\cdot B/Y]/\E[Y^{q/2}]=B\,\E[Y^{q/2-1}]/\E[Y^{q/2}]$.  Let
$h(q)=\ln\E[Y^{q/2}]=\ln\E[e^{q\cdot\frac12\ln Y}]$, the cumulant generating function of $\tfrac12\ln Y$; it is convex (H\"older).  Then
$\langle w\rangle_q=B\exp[h(q-2)-h(q)]$.  For a convex function the secant slope $[h(q)-h(q-2)]/2$ is non-decreasing in $q$, so
$\langle w\rangle_q$ is non-increasing, and strictly decreasing unless $h$ is linear, i.e.\ unless $\ln Y$ is a.s.\ constant at fixed size.
With $c_q=0$ (and $b$ the same for all $q$, as it is the slope of $B$),
\begin{equation}
\frac{\zeta_q}{\zeta_1}=q\,\frac{\langle w\rangle_q}{\langle w\rangle_1}\le q\qquad(q\ge1),
\end{equation}
with equality iff the floor is the same for all firms of equal size.  $\epsilon$ cancels in the ratio.  \exact

\begin{flag}
The abstract and Sec.~VI say the curve ``bends below $q$ only if the floor differs between firms of equal size''.  That is
exact within the class $v^2=(u+B(\Sb))\epsilon^2$ with $B$ a function of size.  Outside it, any dispersion of $\ln v$ that grows
with size bends the curve (Sec.~\ref{sec:cumulant}), and a dispersion that shrinks with size pushes it above $q$.  The
``bin floor $+$ own noise'' counterfactual is exactly this: the per-life $B_i$ is heterogeneous, so the identity does not apply,
and its dispersion is largest among the smallest firms, so by the cumulant criterion the ratios exceed $q$ ($2.08$, $3.65$,
$6.58$).  Stating the scope (``for a law in which the size-dependent part is common to firms of equal size'') would
pre-empt a referee.  The SM paragraph on the composition term already says heterogeneity is necessary in general and
sufficient when $c_q=0$; the main text could carry the same qualifier.
\end{flag}

\subsection{The lognormal case}

If $\ln Y\mid\Sb\sim\mathcal N(m,s^2)$, then $\E[Y^a]=e^{am+a^2s^2/2}$ and
\begin{equation}
\langle w\rangle_q=B\exp\Big[\Big(\frac q2-1\Big)m+\Big(\frac q2-1\Big)^2\frac{s^2}2-\frac q2m-\Big(\frac q2\Big)^2\frac{s^2}2\Big]
=B\,e^{-m}\,e^{(1-q)s^2/2},
\end{equation}
since $(q/2-1)^2-(q/2)^2=1-q$.  Hence
\begin{equation}
\frac{\langle w\rangle_q}{\langle w\rangle_1}=e^{-(q-1)s^2/2},\qquad \frac{\zeta_q}{\zeta_1}=q\,e^{-(q-1)s^2/2}\quad(c_q=0),
\qquad\frac{\langle w\rangle_4}{\langle w\rangle_1}=e^{-3s^2/2}.
\end{equation}
Exact for a lognormal $Y$; but $Y=u+B\ge B$, so $Y$ cannot be exactly lognormal and the formula is an approximation, good
when $s$ is small or $B$ is small compared with the typical $Y$.  \controlled

\paragraph{Narrow floor.}  With $u=\bar ue^{\xi}$ and $\xi$ small, $\ln(u+B)\simeq\ln(\bar u+B)+\frac{\bar u}{\bar u+B}\xi$, so
$s^2\simeq(1-w)^2\Var(\ln u\mid\Sb)$ with $w$ at $\bar u$.  A heterogeneous floor bends the curve only where it is not small against
$B$: at large sizes.  \controlled

\begin{verify}
Monte Carlo with $s^2=0.25$, $10^6$ samples: $\langle w\rangle_q/\langle w\rangle_1=0.8825$, $0.7789$, $0.6876$ against $e^{-(q-1)s^2/2}=0.8825$, $0.7788$,
$0.6873$.  Against SM Table~S6 (power-law graph): $e^{-3s^2/2}$ with the measured $s^2$ gives $0.97$, $0.93$, $0.89$, $0.83$,
$0.68$ against $0.96$, $0.92$, $0.86$, $0.79$, $0.62$ measured, within $0.06$ as stated.
\end{verify}

\subsection{The fitted floor}

The fit $v^2=A\kappa^g+B_j$ makes the floor a deterministic function of the degree and the noise a function of the size bin.
Evaluated with every firm at the mean degree, the floor is common to all firms of a bin, so by the homogeneous-floor result
$\zeta_q=q\zeta_1$ exactly: the $2,3,4$ quoted in the SM is a consequence of the identity, not an independent finding.  With each
firm's own degree, the bend comes from the spread of $\kappa$ within bins.  $g\simeq1$ is what $f\propto\sqrt\kappa$ gives for a floor
proportional to the field's temporal variance.

%======================================================================
\section{Softer self-limitation: the exponent--range trade-off}
\label{sec:theta}

Replace $x_i/m$ by $(x_i/m)^\theta$ in the self-limitation and $x_j/m$ by $(x_j/m)^\theta$ in the interactions.  Without noise,
in relative units $\dot S_i=S_i[1-S_i^\theta+\sum_jJ_{ij}S_j^\theta-\fb]$.  Let $u_i=S_i^\theta$:
\begin{equation}
\dot u_i=\theta S_i^{\theta-1}\dot S_i=\theta u_i\Big[1-\fb-u_i+\sum_jJ_{ij}u_j\Big].
\end{equation}
This is the $\theta=1$ equation in $u$ with time multiplied by $\theta$, \emph{but with the common mode $\fb$ fixed by
$\sum_iS_i=\sum_iu_i^{1/\theta}=N$} rather than by $\sum_iu_i=N$.

\begin{flag}
So $u$ is not exactly the $\theta=1$ model: it is the $\theta=1$ model with a different growth rate.  If $\fb$ self-averages to
$f^*_\theta$, the mapping of Sec.~\ref{sec:relative} applies with $r_\theta=1-f^*_\theta$: $u=(r_\theta/r_1)\,u^{(1)}$ with time rescaled by
$\theta r_\theta/r_1$ rather than $\theta$.  The constant factor on $u$ is harmless (slopes in logs do not see it), but the time factor
changes the effective sampling step, and SM Sec.~S6 shows that the exponents do depend on the step (at $n_{\min}=5$, $\zeta_1$
goes from $0.21$ to $0.16$ between steps $0.5$ and $2$).  Eq.~(13) of the article therefore holds exactly for the
equivalence of the two models at matched steps $\Delta t_u=\theta(r_\theta/r_1)\Delta t_S$, and only approximately at equal steps.  Since
the article uses Eq.~(13) as the noise-free baseline and then shows numerically that noise breaks it, the conclusion is
unaffected; the caveat belongs next to ``exactly if a life's size is its geometric mean''.
\end{flag}

\paragraph{Exponents and ranges.}  $\ln S=\theta^{-1}\ln u$, so every growth rate and every volatility in $S$ is $1/\theta$ times
the one in $u$: $v_S=v_u/\theta$.  If a life's size is its geometric mean, $\ln\Sb=\theta^{-1}\ln\bar u$.  If $v_u\propto\bar u^{-\zeta^{(u)}}$
then $v_S\propto(\Sb^\theta)^{-\zeta^{(u)}}$, so
\begin{equation}
\zeta_1^{(S)}=\theta\,\zeta_1^{(u)},\qquad R_S=\frac{R_u}{\theta},\qquad \zeta_1^{(S)}R_S=\zeta_1^{(u)}R_u .
\end{equation}
With the arithmetic mean the second step is approximate (Jensen).  \exact{} (under the stated conditions).

\paragraph{The noise does not transform.}  With $\dd S=\dots+\sqrt{DS}\,\dd W$, It\^o gives in $u=S^\theta$ the noise term
$\theta S^{\theta-1}\sqrt{DS}\,\dd W=\theta\sqrt D\,S^{\theta-1/2}\dd W=\theta\sqrt D\,u^{1-1/(2\theta)}\dd W$, since $S^{\theta-1/2}=u^{(\theta-1/2)/\theta}$.
At $\theta=1/2$ the exponent is $0$: additive noise.  (The It\^o drift $\tfrac12\theta(\theta-1)S^{\theta-2}DS$ is also added.)  So with
noise the $\theta$ models are genuinely different, and everything after this point in Sec.~VIII is \numer.

\paragraph{$\zeta_1R$ as a log fall.}  If $\ln v$ were linear in $\log_{10}\Sb$ over $R$ decades with slope $-\zeta_1$, the volatility
would fall by $10^{\zeta_1R}$ across the range: $\zeta_1R=\log_{10}(v_{\rm small}/v_{\rm large})$.  At the reference point
$10^{0.40}=2.5$, close to the measured factor $1.11/0.41=2.7$ (the difference is the curvature of the fit).  Compustat:
$10^{0.8}=6.3$ against the factor $5.5$ quoted.  \exact{} (arithmetic).

%======================================================================
\section{Summary of the flags}
\label{sec:flags}

Nothing below invalidates a result of the article.  In order of importance:

\begin{enumerate}
\item \textbf{Scope of ``bends only if the floor differs''} (Sec.~\ref{sec:identity}).  Exact for $v^2=(u+B(\Sb))\epsilon^2$ with $B$
  common to firms of equal size; outside that class the general statement is the cumulant criterion of
  Sec.~\ref{sec:cumulant}.  The main text should carry the qualifier, as the SM does.
\item \textbf{The $\theta$ mapping} (Sec.~\ref{sec:theta}).  The normalisation $\sum S=N$ is not $\sum u=N$, so the time rescaling is
  $\theta r_\theta/r_1$, not $\theta$, and the exponents depend on the step.  Add to the caveat of Eq.~(13).
\item \textbf{Accuracy of the plateau OU formula} (Sec.~\ref{sec:constant}).  The variance error is about $2D/\Phi^2$, $20\%$ at $\Phi=1$
  and $D=0.1$; ``accurate for $\Phi\gtrsim1$'' should say ``to $10\%$ in standard deviation''.
\item \textbf{Validity conditions of the volatility law} (Sec.~\ref{sec:law}).  The linearisation needs $S\gg\sqrt{D/2}$; worth one
  sentence, since it explains the 9\% in the lowest decile.
\item \textbf{The lognormal formula} is exact for lognormal $Y$, but $Y\ge B$; the ``$\simeq$'' in Eq.~(12) is right and could be
  explained in a clause.
\item \textbf{Newborn size in the GLV mapping} (Sec.~\ref{sec:relative}).  In the standard-model variables the entry size $D/r$ and
  the noise $D/r^2$ are two parameters.
\item \textbf{SM Eq.~(S1)}: the $O(N^{-1})\,\dd t$ remainder is exactly zero.
\item \textbf{A sharper statement available}: at $\Phi=0$, $T(D)=\ln(1/D)+2-\gamma_E+O(D)$ (Sec.~\ref{sec:constant}), and for
  $-\Phi\gg\sqrt D$ the $D\to0$ limit is $|\Phi|^{-1}\int_0^{2|\Phi|}e^tE_1(t)\,\dd t$.  Either can replace ``grows as $\ln(1/D)$'' and
  ``nearly independent of $D$''.
\end{enumerate}

\paragraph{What was checked independently for this note.}  The truncated-moment identities; the exact rank-one $\sigma_c$
against Eq.~(6) for $\mu\in\{0,\pm1,-1.76\}$ and two variances; the zero-growth shift $c(\sqrt2)=\tfrac54-\tfrac\pi2$ and slope $-\pi/\sqrt2$;
every quadrature in SM Table~S1 ($P_{\rm est}$ and $T$, all seven rows reproduced to the quoted digits); the Green's function
identity; the two exit-time asymptotics above; the plateau variance by simulation; the lognormal formula by Monte Carlo and
against SM Table~S6; the Slepian dip area.  Everything tagged \numer{} (the one-firm reduction, the field-law exponents, the
per-life decomposition, robustness and the $\theta$ runs) rests on the notebooks and was not re-run.

\bibliographystyle{unsrt}
\bibliography{bib}

\end{document}
